\HydraSection[
  focus x=0.62,
  focus y=0.65,
  radius=0.14,
  zoom=1.85
]{Growth, wounds \& long-range dynamics}


% \HydraSubsection[
%   \mbox{Patterns are stable nonconstant}
%   \mbox{steady states of the full}
%   \mbox{nonlinear reaction--diffusion system}
% ]{What is a pattern?}


\begin{frame}{What do we mean by a pattern?}

  \vspace{-0.34cm}

  \begin{columns}[T,onlytextwidth]

    \begin{column}{0.48\textwidth}
      {\color{HydraMuted}\fontsize{7.5}{9.2}\selectfont
        Return to the Gierer--Meinhardt system
      \par}

      \vspace{-0.30cm}

      {\color{HydraWhite}\fontsize{9.0}{11.5}\selectfont
      \begin{align*}
        \partial_tu
        &=
        D_u\Delta u
        +
        \frac{au^2}{K+v}
        -
        \mu_u u
        \\
        \partial_tv
        &=
        D_v\Delta v
        +
        bu^2
        -
        \mu_v v
      \end{align*}
      }

      \vspace{-0.42cm}

      {\color{HydraMuted}\fontsize{7.3}{9}\selectfont
        with homogeneous Neumann boundary conditions on $\Omega$
      \par}
    \end{column}

    \begin{column}{0.48\textwidth}
      {\centering\color{HydraGold}\bfseries\fontsize{7.5}{9.2}\selectfont
        STEADY STATE\par}

      \vspace{-0.30cm}

      {\color{HydraWhite}\fontsize{8.2}{10.4}\selectfont
      \begin{align*}
        0
        &=
        D_u\Delta\bar u
        +
        \frac{a\bar u^2}{K+\bar v}
        -
        \mu_u\bar u
        \\
        0
        &=
        D_v\Delta\bar v
        +
        b\bar u^2
        -
        \mu_v\bar v
      \end{align*}
      }

      \vspace{-0.42cm}

      \HydraPoint[HydraCyan]
        {Spatial pattern}
        {A stable nonconstant steady state $(\bar u(x),\bar v(x))$}
    \end{column}

  \end{columns}

  \vspace{0.22cm}

  \HydraTakeaway{
    For nonlinear reaction--diffusion systems, a complete characterization of all steady states is generally not known
  }

\end{frame}


\HydraSubsection[
  \mbox{Growth and division change both}
  \mbox{the geometry and the initial data}
  \mbox{of the nonlinear problem}
]{Growth and \mbox{division}}


\begin{frame}{Growth makes the domain part of the dynamics}

  \vspace{-0.12cm}

  \begin{columns}[c,onlytextwidth]
    \begin{column}{0.46\textwidth}
      {\color{HydraGold}\bfseries\fontsize{7.6}{9.1}\selectfont
        TIME-DEPENDENT DOMAIN\par}

      \vspace{-0.16cm}

      {\color{HydraWhite}\bfseries\fontsize{12}{15.2}\selectfont
      \begin{align*}
        \Omega(t)&=[0,L(t)]
        \\
        L'(t)&>0
      \end{align*}
      }

      \vspace{-0.38cm}

      \HydraPoint[HydraCyan]
        {Changing geometry}
        {Growth changes physical distances and the spatial structures that fit inside the domain}

      \vspace{-0.08cm}

      \HydraPoint[HydraGold]
        {Nonlinear question}
        {Does the current pattern persist, reorganize, or generate additional peaks}
    \end{column}

    \begin{column}{0.50\textwidth}
      \centering

      % REPLACEABLE VISUAL
      % Replace the tikzpicture with:
      % \includegraphics[width=\linewidth,height=4.95cm,keepaspectratio]{figures/simulations/growing_domain.png}
      \begin{tikzpicture}[x=0.90cm,y=0.90cm]
        \draw[HydraLine,line width=1.15pt,rounded corners=3pt]
          (0.45,3.30) rectangle (4.05,4.10);
        \fill[HydraGold,opacity=0.65] (1.05,3.34) rectangle (1.35,4.06);

        \draw[-{Latex[length=2.6mm]},HydraWhite,line width=1.0pt]
          (4.35,3.70) -- (5.30,3.70);
        \node[text=HydraMuted,font=\fontsize{6.8}{8.2}\selectfont]
          at (4.82,3.25) {$L(t)$ increases};

        \draw[HydraLine,line width=1.15pt,rounded corners=3pt]
          (0.45,1.25) rectangle (7.65,2.05);
        \fill[HydraGold,opacity=0.65] (1.65,1.29) rectangle (2.00,2.01);
        \fill[HydraGold,opacity=0.65] (5.95,1.29) rectangle (6.30,2.01);

        \draw[<->,HydraCyan,line width=0.9pt]
          (0.45,0.75) -- (7.65,0.75);
        \node[text=HydraCyan,font=\bfseries\fontsize{7.2}{8.6}\selectfont]
          at (4.05,0.35) {new physical scale};
      \end{tikzpicture}
    \end{column}
  \end{columns}

  \vspace{-0.10cm}

  \HydraTakeaway{
    The concentrations evolve while the domain supporting them changes
  }
\end{frame}


% \begin{frame}{The growth experiment begins from a one-peak pattern}

%   \vspace{-0.12cm}

%   \begin{columns}[c,onlytextwidth]
%     \begin{column}{0.70\textwidth}
%       \centering

%       % REPLACEABLE VISUAL
%       % Replace the tikzpicture with:
%       % \includegraphics[width=\linewidth,height=5.25cm,keepaspectratio]{figures/simulations/short_domain_single_peak.png}
%       \begin{tikzpicture}[x=0.90cm,y=0.90cm]
%         \draw[-{Latex[length=2mm]},HydraMuted,line width=0.7pt]
%           (0.45,0.75) -- (8.75,0.75)
%           node[below,text=HydraMuted,font=\fontsize{7}{8.2}\selectfont] {$x$};
%         \draw[HydraGold,line width=1.70pt]
%           plot[smooth] coordinates {(0.55,0.80) (2.10,0.85) (3.25,1.08) (3.85,2.10) (4.20,4.85) (4.58,2.10) (5.15,1.08) (6.30,0.85) (8.60,0.80)};
%         \draw[HydraCyan,line width=1.45pt]
%           plot[smooth] coordinates {(0.55,1.05) (1.70,1.35) (2.80,1.85) (4.20,2.45) (5.65,1.85) (6.85,1.35) (8.60,1.05)};
%         \node[text=HydraGold,font=\bfseries\fontsize{7.2}{8.6}\selectfont]
%           at (5.15,4.55) {$\bar u_1(x)$};
%         \node[text=HydraCyan,font=\bfseries\fontsize{7.2}{8.6}\selectfont]
%           at (6.80,2.15) {$\bar v_1(x)$};
%       \end{tikzpicture}
%     \end{column}

%     \begin{column}{0.26\textwidth}
%       \HydraPoint[HydraGold]
%         {Reference pattern}
%         {For the chosen parameters the initial domain supports a stable one-peak steady state}

%       \vspace{-0.18cm}

%       \HydraPoint[HydraCyan]
%         {Initial condition}
%         {The growth simulation starts from $(\bar u_1(x),\bar v_1(x))$}
%     \end{column}
%   \end{columns}
% \end{frame}


\begin{frame}{Growth can transform one stable pattern into another}

  \vspace{-0.14cm}

  \centering

  % REPLACEABLE VISUAL
  % Replace the tikzpicture with:
  % \includegraphics[width=0.94\linewidth,height=5.30cm,keepaspectratio]{figures/simulations/growth_sequence.png}
  \begin{tikzpicture}[x=0.82cm,y=0.82cm]
    \foreach \shift/\width/\stage in {0/3.2/{$t_0$},4.0/4.7/{$t_1$},9.3/6.0/{$t_2$}}{
      \begin{scope}[shift={(\shift,0)}]
        \draw[HydraLine,line width=1.0pt,rounded corners=2pt]
          (0.20,0.65) rectangle (\width,1.25);
        \node[text=HydraMuted,font=\fontsize{7}{8.4}\selectfont]
          at ({0.5*\width},0.20) {\stage};
      \end{scope}
    }

    \fill[HydraGold,opacity=0.72] (1.38,0.69) rectangle (1.68,1.21);
    \fill[HydraGold,opacity=0.72] (4.75,0.69) rectangle (5.05,1.21);
    \fill[HydraGold,opacity=0.72] (7.65,0.69) rectangle (7.95,1.21);
    \fill[HydraGold,opacity=0.72] (10.20,0.69) rectangle (10.50,1.21);
    \fill[HydraGold,opacity=0.72] (13.65,0.69) rectangle (13.95,1.21);

    \draw[-{Latex[length=2.4mm]},HydraWhite,line width=1.0pt]
      (3.35,0.95) -- (3.85,0.95);
    \draw[-{Latex[length=2.4mm]},HydraWhite,line width=1.0pt]
      (8.85,0.95) -- (9.20,0.95);

    \node[text=HydraGold,font=\bfseries\fontsize{7.3}{8.7}\selectfont]
      at (7.65,3.95) {a new spatial structure can form as the domain grows};
    \node[text=HydraMuted,font=\fontsize{7}{8.5}\selectfont,align=center,text width=9.8cm]
      at (7.65,2.85) {the simulation follows the full nonlinear solution in the physical coordinate};
  \end{tikzpicture}

  \vspace{-0.12cm}

  \HydraTakeaway{
    Growth can reorganize an existing pattern instead of merely stretching it
  }
\end{frame}


\begin{frame}{Division creates two new initial-value problems}

  \vspace{-0.12cm}

  \begin{columns}[c,onlytextwidth]
    \begin{column}{0.67\textwidth}
      \centering

      % REPLACEABLE VISUAL
      % Replace the tikzpicture with:
      % \includegraphics[width=\linewidth,height=5.10cm,keepaspectratio]{figures/simulations/domain_division.png}
      \begin{tikzpicture}[x=0.88cm,y=0.88cm]
        \draw[HydraLine,line width=1.15pt,rounded corners=3pt]
          (0.45,3.50) rectangle (8.55,4.30);
        \fill[HydraGold,opacity=0.72] (1.10,3.54) rectangle (1.45,4.26);
        \fill[HydraGold,opacity=0.52] (7.40,3.54) rectangle (7.68,4.26);
        \draw[HydraRed,dashed,line width=1.2pt]
          (4.55,3.25) -- (4.55,4.60);
        \node[text=HydraRed,font=\bfseries\fontsize{7.2}{8.6}\selectfont]
          at (4.55,4.95) {cut};

        \draw[-{Latex[length=2.5mm]},HydraWhite,line width=1.0pt]
          (4.55,2.95) -- (2.70,2.10);
        \draw[-{Latex[length=2.5mm]},HydraWhite,line width=1.0pt]
          (4.55,2.95) -- (6.45,2.10);

        \draw[HydraLine,line width=1.10pt,rounded corners=3pt]
          (0.55,0.70) rectangle (3.85,1.40);
        \fill[HydraGold,opacity=0.72] (1.10,0.74) rectangle (1.45,1.36);

        \draw[HydraLine,line width=1.10pt,rounded corners=3pt]
          (5.20,0.70) rectangle (8.50,1.40);
        \fill[HydraGold,opacity=0.52] (7.40,0.74) rectangle (7.68,1.36);

        \node[text=HydraMuted,font=\fontsize{6.8}{8.2}\selectfont]
          at (2.20,0.25) {fragment A};
        \node[text=HydraMuted,font=\fontsize{6.8}{8.2}\selectfont]
          at (6.85,0.25) {fragment B};
      \end{tikzpicture}
    \end{column}

    \begin{column}{0.29\textwidth}
      \HydraPoint[HydraGold]
        {Inherited profiles}
        {Each fragment receives the restriction of the pre-cut profiles of $u$ and $v$}

      \vspace{-0.18cm}

      \HydraPoint[HydraCyan]
        {New boundaries}
        {No-flux conditions are imposed at the newly created ends}

      \vspace{-0.18cm}

      \HydraPoint[HydraRed]
        {Independent evolution}
        {After division the fragments solve two separate nonlinear problems}
    \end{column}
  \end{columns}
\end{frame}


% \begin{frame}{The fragments inherit different initial profiles}

%   \vspace{-0.12cm}

%   \begin{columns}[c,onlytextwidth]
%     \begin{column}{0.64\textwidth}
%       \centering

%       % REPLACEABLE VISUAL
%       % Replace the tikzpicture with:
%       % \includegraphics[width=\linewidth,height=5.15cm,keepaspectratio]{figures/simulations/fragment_fates.png}
%       \begin{tikzpicture}[x=0.88cm,y=0.88cm]
%         \draw[-{Latex[length=1.9mm]},HydraMuted,line width=0.7pt]
%           (0.35,3.40) -- (7.65,3.40);
%         \draw[HydraGold,line width=1.55pt]
%           plot[smooth] coordinates {(0.45,3.45) (1.15,3.50) (1.55,4.05) (1.78,5.35) (2.02,4.00) (2.50,3.50) (7.50,3.45)};
%         \node[text=HydraGold,font=\bfseries\fontsize{7.2}{8.6}\selectfont]
%           at (5.45,4.65) {fragment containing a peak};

%         \draw[-{Latex[length=1.9mm]},HydraMuted,line width=0.7pt]
%           (0.35,0.85) -- (7.65,0.85);
%         \draw[HydraCyan,line width=1.55pt]
%           plot[smooth] coordinates {(0.45,0.90) (1.60,1.15) (2.80,1.40) (4.00,1.25) (5.30,1.05) (7.50,0.90)};
%         \node[text=HydraCyan,font=\bfseries\fontsize{7.2}{8.6}\selectfont]
%           at (5.25,1.85) {fragment without a mature peak};

%         \draw[HydraRed,dashed,line width=0.85pt]
%           (0.45,2.25) -- (7.50,2.25);
%         \node[anchor=west,text=HydraRed,font=\fontsize{6.8}{8.2}\selectfont]
%           at (0.55,2.45) {schematic activation threshold};
%       \end{tikzpicture}
%     \end{column}

%     \begin{column}{0.32\textwidth}
%       \HydraPoint[HydraGold]
%         {Peak-containing fragment}
%         {A sufficiently strong inherited organizer can remain self-sustaining}

%       \vspace{-0.18cm}

%       \HydraPoint[HydraCyan]
%         {Headless fragment}
%         {A weak residual profile can remain in the basin of the zero steady state}

%       \vspace{-0.18cm}

%       \HydraPoint[HydraRed]
%         {Different initial data}
%         {The cut changes both the geometry and the position relative to the nonlinear threshold}
%     \end{column}
%   \end{columns}
% \end{frame}


\HydraSubsection[
  \mbox{Regeneration requires a localized}
  \mbox{wound signal strong enough to cross}
  \mbox{the nonlinear activation threshold}
]{Role of the wound}


\begin{frame}{A cut alone does not guarantee regeneration}

  \vspace{-0.12cm}

  \begin{columns}[c,onlytextwidth]
    \begin{column}{0.70\textwidth}
      \centering

      % REPLACEABLE VISUAL
      % Replace the tikzpicture with:
      % \includegraphics[width=\linewidth,height=5.10cm,keepaspectratio]{figures/simulations/no_wound_decay.png}
      \begin{tikzpicture}[x=0.76cm,y=0.82cm]
        \foreach \shift/\amp/\stage in {0/2.6/{$t_0$},3.8/1.5/{$t_1$},7.6/0.5/{$t_2$}}{
          \begin{scope}[xshift=\shift cm]
            \draw[-{Latex[length=1.8mm]},HydraMuted,line width=0.65pt]
              (0.25,0.70) -- (3.55,0.70);
            \draw[HydraCyan,line width=1.45pt]
              plot[smooth] coordinates {(0.35,0.75) (0.85,0.80) (1.35,{0.75+0.35*\amp}) (1.75,{0.75+\amp}) (2.15,{0.75+0.35*\amp}) (2.65,0.80) (3.45,0.75)};
            \node[text=HydraMuted,font=\fontsize{6.8}{8.2}\selectfont]
              at (1.85,0.25) {\stage};
          \end{scope}
        }

        \draw[-{Latex[length=2.3mm]},HydraWhite,line width=0.95pt]
          (3.45,1.75) -- (3.70,1.75);
        \draw[-{Latex[length=2.3mm]},HydraWhite,line width=0.95pt]
          (7.25,1.75) -- (7.50,1.75);
        \node[text=HydraRed,font=\bfseries\fontsize{8}{9.6}\selectfont]
          at (5.65,4.65) {$s_{\mathrm w}\equiv0$};
      \end{tikzpicture}
    \end{column}

    \begin{column}{0.26\textwidth}
      \HydraPoint[HydraRed]
        {Simulation outcome}
        {Without wound input the headless fragment can converge to $(0,0)$}

      \vspace{-0.18cm}

      \HydraPoint[HydraCyan]
        {Nonlinear explanation}
        {The inherited profile remains below the activation threshold and inside the basin of $(0,0)$}
    \end{column}
  \end{columns}

  \vspace{-0.12cm}

  \HydraTakeaway{
    A new boundary does not by itself move the system into the basin of a nonzero pattern
  }
\end{frame}


\begin{frame}{A wound provides a localized transient signal}

  \vspace{-0.12cm}

  \begin{columns}[c,onlytextwidth]
    \begin{column}{0.56\textwidth}
      {\color{HydraGold}\bfseries\fontsize{7.6}{9.1}\selectfont
        ACTIVATOR EQUATION WITH WOUND INPUT\par}

      \vspace{-0.26cm}

      {\color{HydraWhite}\fontsize{9.7}{12.4}\selectfont
      \begin{align*}
        \partial_tu
        &=D_u\Delta u
        +\frac{au^2}{K+v}
        -\mu_u u
        +s_{\mathrm w}(x,t)
      \end{align*}
      }

      \vspace{-0.44cm}

      {\color{HydraWhite}\fontsize{8.2}{10.4}\selectfont
      \begin{align*}
        \operatorname{supp}s_{\mathrm w}
        &\subset
        [0,\ell_{\mathrm w}]
        \times
        [0,T_{\mathrm w}]
      \end{align*}
      }

      \vspace{-0.42cm}

      \HydraPoint[HydraGold]
        {Localized in space}
        {The perturbation acts near the newly created wound boundary}

      \vspace{-0.18cm}

      \HydraPoint[HydraCyan]
        {Transient in time}
        {The input disappears after a finite activation period}
    \end{column}

    \begin{column}{0.40\textwidth}
      \centering

      % REPLACEABLE VISUAL
      % Replace the tikzpicture with:
      % \includegraphics[width=\linewidth,height=4.90cm,keepaspectratio]{figures/simulations/wound_signal.png}
      \begin{tikzpicture}[x=0.88cm,y=0.88cm]
        \draw[-{Latex[length=2mm]},HydraMuted,line width=0.7pt]
          (0.35,0.70) -- (5.45,0.70)
          node[below,text=HydraMuted,font=\fontsize{7}{8.2}\selectfont] {$x$};
        \draw[-{Latex[length=2mm]},HydraMuted,line width=0.7pt]
          (0.50,0.50) -- (0.50,4.95)
          node[left,text=HydraMuted,font=\fontsize{7}{8.2}\selectfont] {$s_{\mathrm w}$};

        \fill[HydraGold,opacity=0.13] (0.50,0.70) rectangle (1.55,4.55);
        \draw[HydraGold,line width=1.65pt]
          plot[smooth] coordinates {(0.55,4.30) (0.80,4.20) (1.05,3.75) (1.35,2.10) (1.70,1.10) (2.25,0.78) (5.30,0.72)};
        \node[anchor=west,text=HydraGold,font=\bfseries\fontsize{7}{8.4}\selectfont]
          at (1.65,4.10) {wound boundary};

        \draw[<->,HydraCyan,line width=0.9pt]
          (3.25,1.20) -- (3.25,3.05);
        \node[anchor=west,text=HydraCyan,font=\fontsize{6.8}{8.2}\selectfont]
          at (3.42,2.12) {pulse amplitude};
      \end{tikzpicture}

      \vspace{-0.18cm}

      \HydraPoint[HydraRed]
        {Threshold crossing}
        {Sufficient amplitude and duration move the local state across the activation threshold}
    \end{column}
  \end{columns}
\end{frame}


\begin{frame}{A transient wound can create a persistent organizer}

  \vspace{-0.14cm}

  \centering

  % REPLACEABLE VISUAL
  % Replace the tikzpicture with:
  % \includegraphics[width=0.94\linewidth,height=5.25cm,keepaspectratio]{figures/simulations/wound_regeneration_sequence.png}
  \begin{tikzpicture}[x=0.78cm,y=0.80cm]
    \foreach \shift/\stage in {0/{cut},3.9/{wound pulse},7.8/{peak initiation},11.7/{self-maintenance}}{
      \begin{scope}[xshift=\shift cm]
        \draw[HydraLine,line width=1.0pt,rounded corners=2pt]
          (0.25,0.70) rectangle (3.25,1.30);
        \node[text=HydraMuted,font=\fontsize{6.6}{8}\selectfont,align=center,text width=2.5cm]
          at (1.75,0.20) {\stage};
      \end{scope}
    }

    \draw[HydraRed,line width=1.5pt] (0.55,0.58) -- (0.55,1.42);
    \fill[HydraGold,opacity=0.35] (4.15,0.74) rectangle (5.10,1.26);
    \fill[HydraGold,opacity=0.58] (8.05,0.74) rectangle (8.42,1.26);
    \fill[HydraGold,opacity=0.80] (11.95,0.74) rectangle (12.32,1.26);

    \draw[-{Latex[length=2.2mm]},HydraWhite,line width=0.9pt] (3.25,1.00) -- (3.80,1.00);
    \draw[-{Latex[length=2.2mm]},HydraWhite,line width=0.9pt] (7.15,1.00) -- (7.70,1.00);
    \draw[-{Latex[length=2.2mm]},HydraWhite,line width=0.9pt] (11.05,1.00) -- (11.60,1.00);

    \node[text=HydraGold,font=\bfseries\fontsize{7.5}{9}\selectfont]
      at (7.60,4.30) {the wound input disappears while the newly initiated peak persists};
    \node[text=HydraMuted,font=\fontsize{7}{8.5}\selectfont,align=center,text width=10.5cm]
      at (7.60,3.15) {after threshold crossing, the autonomous GM dynamics can approach a stable nonconstant steady state};
  \end{tikzpicture}

  \vspace{-0.12cm}

  \HydraTakeaway{
    The wound initiates regeneration while the nonlinear feedback maintains the resulting organizer
  }
\end{frame}


\begin{frame}{The simulations expose a missing biological mechanism}

  \vspace{-0.12cm}

  \begin{columns}[c,onlytextwidth]
    \begin{column}{0.47\textwidth}
      {\color{HydraGold}\bfseries\fontsize{7.7}{9.2}\selectfont
        WHAT THE SIMULATIONS DEMONSTRATE\par}

      \vspace{0.10cm}

      \HydraPoint[HydraGold]
        {Several stable patterns}
        {The full nonlinear system can support different steady spatial profiles}

      \vspace{-0.18cm}

      \HydraPoint[HydraCyan]
        {Geometry-dependent outcomes}
        {Growth and division change the domain and the initial-value problem}

      \vspace{-0.18cm}

      \HydraPoint[HydraViolet]
        {Wound-dependent regeneration}
        {A transient localized input can initiate a persistent organizer}
    \end{column}

    \begin{column}{0.49\textwidth}
      {\color{HydraRed}\bfseries\fontsize{7.7}{9.2}\selectfont
        WHAT THE MODEL STILL ASSUMES\par}

      \vspace{0.10cm}

      \HydraPoint[HydraRed]
        {Long-range molecular inhibition}
        {The classical mechanism requires an inhibitor that spreads farther than the activator}

      \vspace{-0.18cm}

      \HydraPoint[HydraMuted]
        {Uncertain biological identity}
        {Candidate activators exist, but the corresponding rapidly spreading inhibitor remains unclear}

      \vspace{0.12cm}

      \begin{center}
        {\color{HydraWhite}\bfseries\fontsize{11.2}{14}\selectfont
          Could mechanics supply the missing inhibition?
        \par}
      \end{center}
    \end{column}
  \end{columns}

  \vspace{-0.10cm}

  \HydraTakeaway{
    This motivates a mechanochemical model in which tissue deformation participates in pattern formation
  }
\end{frame}
